\documentclass[11pt]{article}
\usepackage{fontspec,amsmath,amssymb,amsthm,hyperref}
\usepackage[margin=1in]{geometry}
\setlength{\emergencystretch}{3em}
\title{A concrete independent initial phase and canonical Brownian source}
\author{LangevinLean mathematical exposition}
\date{}
\begin{document}\maketitle
\noindent Source module: \texttt{KineticLangevin/CanonicalProduct.lean}. Compile this proof document with XeLaTeX or LuaLaTeX. The accompanying HTML pairs the mathematical argument with the exact Lean statements and records the scope of the formalization.


The stationarity proof takes place on one explicit probability space: the product of an initial phase law and the constructed vector Brownian source. This module builds the actual initial L\ensuremath{{}^{2}} variables, the actual continuous Brownian forcing and its OU-filtered innovations, then proves the complete initial/noise joint law. The resulting splitting law identity has no assumed independence or Gaussian marginal fields.


\section{Represent the original initial phase on the product source}


\texttt{H} is the finite-dimensional Euclidean state, and \texttt{Sample} pairs a phase with a vector Brownian sample. The probability measure is exactly \(\pi\otimes P_B\). Pull the velocity and position projections through the first product projection. Their square integrability follows from that of the original phase law. The two \texttt{toLp} representative identities prove that the resulting phase agrees almost everywhere with the original initial coordinate.



\noindent Lean declarations: \texttt{H}, \texttt{Sample}, \texttt{probability}, \texttt{initial}, \texttt{initial\_coe}.

\section{Lift the actual continuous Brownian input}


The second product projection preserves the Brownian source law. Its pullback is a linear isometry on L\ensuremath{{}^{2}}, hence a continuous linear map. Apply it to \(\sqrt{2\gamma}B_t\), where the constructed Brownian input is continuous in L\ensuremath{{}^{2}} and zero at time zero. Linearity and continuity give the same properties after adjoining the independent initial phase. Thus the forcing used by the deterministic integral solver is an actual Brownian forcing on this very source.



\noindent Lean declarations: \texttt{brownianLift}, \texttt{input}, \texttt{continuous\_input}, \texttt{input\_zero}.

\section{Use the exact OU-filtered Brownian intervals}


The innovation in interval \([nh,(n+1)h]\) is the pulled-back isonormal convolution \(\sqrt{2\gamma}\int_{nh}^{(n+1)h}e^{-\gamma((n+1)h-s)}dB_s\). Its representative is exactly the original source variable composed with the Brownian projection. Continuous linear maps commute with the deterministic OU filter. Combined with the proved canonical integration-by-parts identity, this shows the continuous-input filter is exactly this innovation, not merely a random variable with a matching covariance.



\noindent Lean declarations: \texttt{innovation}, \texttt{innovation\_coe}, \texttt{ouInputIncrement\_input}.

\section{Derive the complete independent Gaussian stream law}


The canonical OU marginal theorem gives covariance \((1-e^{-2\gamma h})I\); the disjoint-interval theorem gives independence of the full vector innovations across time. Their whole stream therefore has the infinite product of the thermal Gaussian law. The explicit initial-law/source product makes the initial phase independent of that entire stream. Finally use the almost-everywhere initial and innovation identities to obtain the complete joint law of the actual L\ensuremath{{}^{2}} variables.



\noindent Lean declarations: \texttt{initial\_innovation\_jointLaw}.

\section{Identify both actual splitting endpoint laws}


Apply the proved full-input-law theorem to the actual lifted Verlet–OU chain. Its phase pushforward is exactly the Markov iterate with fresh thermal Gaussian OU innovations and initial law \(\pi\). The exact OU-filter identity also identifies the continuous-input chain with the same noise-driven chain, so its endpoint has that law as well. These identities provide the missing connection between the pathwise convergence theorem and the measure-level weak stationarity calculation.



\noindent Lean declarations: \texttt{splitNoiseChain\_law}, \texttt{splitChain\_law}.

\end{document}
